% =============================================================================
% Section 8: Growth Prediction — A Three-Stage Complexity Ladder
% =============================================================================
\section{Growth Prediction: A Three-Stage Complexity Ladder}
\label{sec:growth}

% ---------------------------------------------------------------------------
% 8.1 Overview and Design Rationale
% ---------------------------------------------------------------------------
\subsection{Overview and Design Rationale}
\label{sec:growth:overview}

This section presents the core methodological contribution: a three-stage
complexity ladder for meningioma growth prediction, where each stage must
justify its existence through improved out-of-sample performance over the
preceding one. The key design principle is parsimony: at $N = 31$--$58$
patients with $\sim$112 total observations, models with more than 2--3
predictor parameters per patient risk overfitting, and complex architectures
must demonstrate they capture signal rather than noise.

\subsubsection{Why Not a Neural ODE?}

The original pipeline design used a Neural ODE
\cite{chen2018neural,rubanova2019latent} for temporal modelling. This was
abandoned due to catastrophic overparameterisation: with 33~patients
(100~observations, 112~forward pairs), the Neural ODE's $\sim$3,100+
parameters produce a parameter-to-observation ratio of approximately 27.7.
The full pivot rationale is documented in \cite{pascual2026pivot}.

\subsubsection{Stage Summary}

The three stages, in order of increasing complexity:
\begin{description}[leftmargin=*,labelwidth=!]
  \item[Stage 1:] Segment tumour $\to$ extract volume $\to$ model scalar
    volume trajectories. Parameters: 2--5 (population level).
  \item[Stage 2:] Estimate patient-specific latent severity $s_i \in [0, 1]$
    $\to$ model growth as $g(s_i, t; \boldsymbol{\theta})$. Parameters:
    population $\boldsymbol{\theta}$ + 1 per patient.
  \item[Stage 3:] BrainSegFounder features $\to$ PCA $\to$ GP with ARD
    kernel. Parameters: population kernel hyperparameters + 3--6 PCA
    components.
\end{description}

A cross-cutting variance decomposition (\cref{sec:eval:variance}) quantifies
the marginal $\Delta R^2$ of each stage.

% ---------------------------------------------------------------------------
% 8.2 Stage 1: Volumetric Growth Models
% ---------------------------------------------------------------------------
\subsection{Stage 1: Segmentation-Based Volumetric Baseline}
\label{sec:growth:stage1}

Stage~1 is the simplest credible growth prediction approach: segment the
tumour from each MRI, compute whole-tumour volume in mm$^3$, and model the
resulting scalar volume trajectories. Three models of increasing flexibility
are evaluated.

\subsubsection{Model 1A: Scalar GP (ScalarGP)}

A single GP is fitted to the pooled volume observations across all patients:
\begin{equation}
  v_i(t_{ij}) \sim \GP\bigl(m(t),\; k_{5/2}(t, t') + \sigma_n^2 \delta(t, t')\bigr),
  \label{eq:scalargp}
\end{equation}
where $v_i(t_{ij}) = \log(V_{i,\text{WT}}(t_{ij}) + 1)$ is the
log-transformed whole-tumour volume, $m(t) = \beta_0 + \beta_1 t$ is a
linear mean function, and $k_{5/2}$ is the Mat\'{e}rn-5/2 kernel
(\cref{eq:matern52}).

\textbf{Hyperparameters:} $\{\beta_0, \beta_1, \sigma_f, \ell, \sigma_n\}$
= 5 total, estimated by maximising the log-marginal likelihood.

\textbf{Covariates:} Age and sex can be incorporated into the mean function:
\begin{equation}
  m_i(t) = \beta_0 + \beta_1 t + \beta_2 \,\text{age}_i
    + \beta_3 \,\text{sex}_i,
  \label{eq:scalargp_covariates}
\end{equation}
adding 2~parameters. Given the sample size constraints ($N/15 \approx 2$--$3$
allowed parameters), covariates are included only if they improve LOPO-CV
performance.

\subsubsection{Model 1B: Linear Mixed-Effects Model (LME)}

The LME introduces patient-specific random effects to separate population
and individual trends:
\begin{equation}
  v_i(t_{ij}) = \underbrace{(\beta_0 + b_{0i})}_{\text{intercept}}
    + \underbrace{(\beta_1 + b_{1i})}_{\text{slope}} \cdot t_{ij}
    + \epsilon_{ij},
  \label{eq:lme_vol}
\end{equation}
where $(\beta_0, \beta_1)$ are fixed effects (population intercept and
slope), $(b_{0i}, b_{1i}) \sim \mathcal{N}(\mathbf{0}, \boldsymbol{\Omega})$
are patient-specific random effects, and
$\epsilon_{ij} \sim \mathcal{N}(0, \sigma^2)$ is residual noise.

\textbf{Parameter count:} 2~fixed effects $(\beta_0, \beta_1)$, 3~variance
parameters (2~diagonal + 1~off-diagonal of $\boldsymbol{\Omega}$),
1~residual variance $(\sigma^2)$ = 6 total. Estimated via REML using
\texttt{statsmodels.MixedLM}.

\textbf{Single-patient inference:} The Best Linear Unbiased Predictor (BLUP)
\cite{robinson1991blup} provides patient-specific predictions with automatic
shrinkage: patients with few observations ($n_i = 2$) receive predictions
closer to the population mean, while patients with more observations
($n_i \geq 4$) receive data-driven predictions.

\subsubsection{Model 1C: Hierarchical GP (HGP)}

The HGP extends the LME by adding nonlinear residuals via a GP:
\begin{equation}
  v_i(t) \sim \GP\bigl(m(t),\; k_{5/2}(t, t') + \sigma_n^2 \delta(t, t')\bigr),
  \label{eq:hgp_vol}
\end{equation}
where the mean function is the LME population trend:
\begin{equation}
  m(t) = \hat{\beta}_0 + \hat{\beta}_1 \cdot t,
  \label{eq:gp_mean_lme}
\end{equation}
using fixed effects estimated from the LME. The GP captures nonlinear
deviations from the population trend.

\textbf{Hierarchical hyperparameter sharing:} Shared kernel hyperparameters
$\boldsymbol{\theta} = (\sigma_f, \ell, \sigma_n)$ are estimated by
maximising the pooled log-marginal likelihood across all patients:
\begin{equation}
  \hat{\boldsymbol{\theta}} = \arg\max_{\boldsymbol{\theta}}
    \sum_{i=1}^{N} \log p\bigl(v_i \mid
    \mathbf{t}_i, \boldsymbol{\theta}\bigr).
  \label{eq:pooled_marglik}
\end{equation}

\textbf{Parameter count:} 3~kernel hyperparameters (shared) + 2~mean
function parameters (from LME) = 5. Optimised via L-BFGS-B with 5~random
restarts.

\textbf{Model nesting:} The GP with linear mean function $\supset$ LME
(\cref{sec:bg:gp}), ensuring the HGP can never perform worse than the LME
on the training data. Improvement must come from the nonlinear Mat\'{e}rn
component, which we evaluate via log-marginal likelihood comparison.

\subsubsection{Gompertz Mean Function}

As an alternative to the linear population mean, we evaluate a Gompertz
mean function motivated by the growth biology literature
(\cref{sec:bg:growth_biology}):
\begin{equation}
  m_G(t) = V_0 \exp\!\left(\frac{\alpha}{\beta}
    \bigl(1 - e^{-\beta t}\bigr)\right),
  \label{eq:gompertz_mean}
\end{equation}
where $V_0$ is the initial volume, $\alpha$ is the initial growth rate,
and $\beta$ is the deceleration parameter. In log-volume space this becomes
$\log m_G(t) = \log V_0 + \frac{\alpha}{\beta}(1 - e^{-\beta t})$. The
Gompertz mean adds 1~parameter ($\beta$) beyond the linear mean, and is
evaluated as ablation~A2 (\cref{sec:eval:ablations}).

% ---------------------------------------------------------------------------
% 8.3 Stage 2: Latent Severity Model
% ---------------------------------------------------------------------------
\subsection{Stage 2: Latent Severity Model}
\label{sec:growth:stage2}

Stage~2 posits that inter-patient growth heterogeneity is largely explained
by a single latent severity variable $s_i \in [0, 1]$. This is motivated by
Vaghi et al.'s \cite{vaghi2020reduced} finding that a single patient-specific
parameter captures most growth variability in tumour populations.

\subsubsection{Quantile Transform}

To place all patients on a common scale, we apply a quantile transform to
both time and growth:
\begin{enumerate}[leftmargin=*]
  \item \textbf{Temporal normalisation.} Within each patient, times are
    normalised to $[0, 1]$:
    $\tilde{t}_{ij} = (t_{ij} - t_{i,\min}) / (t_{i,\max} - t_{i,\min})$.
  \item \textbf{Growth normalisation.} Volume changes are mapped to
    quantiles of the population distribution:
    $q_{ij} = F_{\text{pop}}(\Delta v_{ij})$, where $F_{\text{pop}}$ is the
    empirical CDF of all observed volume changes across the training set.
\end{enumerate}

The quantile-transformed observations $q_{ij} \in [0, 1]$ represent the
relative position of each observation within the population growth
distribution.

\subsubsection{NLME Formulation}

The severity model is a nonlinear mixed-effects model:
\begin{equation}
  q_{ij} = g(s_i, \tilde{t}_{ij}; \boldsymbol{\theta}) + \epsilon_{ij},
  \quad \epsilon_{ij} \sim \mathcal{N}(0, \sigma^2),
  \label{eq:severity_nlme}
\end{equation}
where $g$ is a monotonic growth function satisfying boundary conditions
$g(s, 0) = 0$ (no growth at baseline) and $g(s, 1) \leq s$ (severity bounds
maximum growth). The latent severity $s_i \sim \text{Beta}(\alpha, \beta)$ has
a population prior.

\subsubsection{Growth Function Options}

Three growth functions are evaluated for $g$:

\begin{description}[leftmargin=*,labelwidth=!]
  \item[Sigmoid.] The simplest parametric option:
    $g(s, t; \boldsymbol{\theta}) = s \cdot \sigma(\theta_1 (t - \theta_2))$,
    where $\sigma$ is the logistic sigmoid and
    $\boldsymbol{\theta} = (\theta_1, \theta_2)$ controls the steepness and
    inflection point. Monotonic in $s$ by construction. 2~population
    parameters.

  \item[Reduced Gompertz.] Following Vaghi et al.\
    \cite{vaghi2020reduced}:
    $g(s, t; \beta) = s \cdot [1 - \exp(-\alpha(s)(1 - e^{-\beta t})/\beta)]$,
    where $\alpha(s)$ maps severity to growth rate. 1~population parameter
    ($\beta$).

  \item[CMNN.] A constrained monotonic neural network
    \cite{runje2023cmnn} with non-negative weights ensuring
    $\partial g / \partial s \geq 0$ everywhere. The most flexible option
    but with the highest parameter count ($\sim$10--20 weights). Evaluated
    only if the parametric options prove insufficient.
\end{description}

These are compared in ablation~A3 (\cref{sec:eval:ablations}).

\subsubsection{Joint Estimation}

The population parameters $\boldsymbol{\theta}$ and patient severities
$\{s_i\}_{i=1}^N$ are estimated jointly by maximising the marginal
likelihood:
\begin{equation}
  \hat{\boldsymbol{\theta}} = \arg\max_{\boldsymbol{\theta}}
    \sum_{i=1}^{N} \log \int p(\mathbf{q}_i \mid s_i, \boldsymbol{\theta})
    \, p(s_i) \, ds_i.
  \label{eq:severity_mle}
\end{equation}
For the Beta prior, the inner integral is computed via Gauss--Hermite
quadrature or Laplace approximation.

\subsubsection{Test-Time Severity Estimation}

For a held-out patient with observed trajectory
$\{(t_{ij}, q_{ij})\}_{j=1}^{n_i}$, the posterior severity is:
\begin{equation}
  p(s_i \mid \mathbf{q}_i, \hat{\boldsymbol{\theta}}) \propto
    p(\mathbf{q}_i \mid s_i, \hat{\boldsymbol{\theta}}) \, p(s_i),
  \label{eq:severity_posterior}
\end{equation}
which is maximised or integrated over for point and distributional predictions
respectively. For patients with $n_i = 2$ (the most common case), the
posterior is dominated by the prior, yielding predictions close to the
population median---a desirable behaviour analogous to LME shrinkage.

\subsubsection{Severity Estimation from Baseline Features}

A secondary objective of Stage~2 is to predict severity from baseline
(single-timepoint) information, enabling predictions without longitudinal
data. Candidate predictors include:
\begin{itemize}[leftmargin=*]
  \item Clinical covariates: age, sex, tumour location.
  \item Baseline volume $V_{i,0}$.
  \item BrainSegFounder features at $t_0$ (linking to Stage~3).
\end{itemize}
This is evaluated in ablation~A4 (\cref{sec:eval:ablations}).

% ---------------------------------------------------------------------------
% 8.4 Stage 3: Representation-Learning Augmented Prediction
% ---------------------------------------------------------------------------
\subsection{Stage 3: Representation-Learning Augmented Prediction}
\label{sec:growth:stage3}

Stage~3 asks whether high-dimensional deep features from BrainSegFounder
capture growth-relevant information beyond what volume and severity already
explain. This stage uses the LoRA-adapted encoder (\cref{sec:lora}) and
SDP (\cref{sec:sdp}) from the original pipeline, but positions them as the
final complexity level rather than the primary approach.

\subsubsection{Feature Extraction and Compression}

The frozen pipeline (encoder + SDP) produces per-scan latent vectors
$\mathbf{z} \in \R^{128}$. For growth prediction, these are compressed via
within-fold PCA to $\tilde{\mathbf{z}} \in \R^{d_{\text{PCA}}}$ with
$d_{\text{PCA}} \in \{3, 4, 5, 6\}$ selected by the sample size constraint
and cross-validated performance. PCA is fitted on the training set within
each LOPO-CV fold to prevent information leakage.

\subsubsection{GP with ARD Kernel}

The compressed features are used as inputs to a GP with an ARD kernel:
\begin{equation}
  k_{\text{ARD}}(\tilde{\mathbf{z}}, \tilde{\mathbf{z}}'; t, t') =
    k_{5/2}(t, t') \cdot \prod_{d=1}^{d_{\text{PCA}}}
    k_{\text{SE}}\!\left(\frac{\tilde{z}^{(d)} - \tilde{z}'^{(d)}}{\ell_d}\right)
    + \sigma_n^2 \delta_{ii'} \delta(t, t'),
  \label{eq:ard_kernel}
\end{equation}
where $\ell_d$ are per-dimension length-scales. Dimensions with large
$\ell_d$ are effectively deactivated, providing automatic feature selection.
The mean function is the Stage~1 population trend.

\textbf{Parameter count:} $d_{\text{PCA}}$ ARD length-scales + 3 temporal
kernel hyperparameters + mean function parameters. For $d_{\text{PCA}} = 4$:
$\sim$9 total hyperparameters, estimated from the pooled marginal likelihood.

\subsubsection{Integration with Earlier Stages}

Stage~3 does not replace Stages~1 and~2 but augments them. The deep features
enter as additional covariates in the GP mean function or as ARD inputs
alongside time, while the population trend from Stage~1 remains the prior
mean. The marginal contribution is quantified by the variance decomposition
(\cref{sec:eval:variance}).

\subsubsection{SDP with Supervised Residual Targets}

An alternative to PCA compression is to use the SDP's supervised partition
$\Vol \in \R^{32}$ directly. The SDP is designed to place volume-relevant
information in a structured subspace. However, 32~dimensions is still far
too many for GP modelling at $N = 33$. Two options are evaluated:
\begin{enumerate}[leftmargin=*]
  \item PCA on $\Vol$ alone (expected to require fewer components than PCA
    on full $\mathbf{z}$, since the SDP has already concentrated
    volume-relevant information).
  \item Using $\Vol$ as input to a GP with an ICM coregionalisation matrix
    (\cref{eq:icm}), treating each latent dimension as an output channel.
    This is the PA-MOGP structure from the original design.
\end{enumerate}

\begin{remark}[Stage~3 is conditional]
  Stage~3 is attempted only if Stages~1 and~2 leave unexplained variance
  that could plausibly be captured by additional features. If the variance
  decomposition shows that volume and severity together explain $R^2 > 0.70$,
  the marginal room for improvement from deep features is limited, and
  Stage~3 serves primarily as a negative result confirming parsimony.
\end{remark}

% ---------------------------------------------------------------------------
% 8.5 LOPO-CV Protocol
% ---------------------------------------------------------------------------
\subsection{Leave-One-Patient-Out Cross-Validation}
\label{sec:growth:lopo}

All three stages share a common LOPO-CV protocol. With only 31--58~patients,
$k$-fold CV leaves too few patients per fold for stable model estimation.
LOPO-CV uses maximum training data per fold:

\begin{algorithm}[H]
  \caption{LOPO-CV Evaluation Protocol (shared across all stages)}
  \label{alg:lopo}
  \begin{algorithmic}[1]
    \FOR{$i = 1$ to $N$}
      \STATE $\mathcal{D}_{\text{train}} \leftarrow
        \{\mathcal{T}_j : j \neq i\}$ \COMMENT{$N-1$ patients}
      \STATE $\mathcal{D}_{\text{test}} \leftarrow \{\mathcal{T}_i\}$
        \COMMENT{held-out patient}
      \STATE \textbf{Within-fold operations:}
      \STATE \quad Fit PCA on $\mathcal{D}_{\text{train}}$ (Stage~3 only)
      \STATE \quad Estimate severity $\{s_j\}_{j \neq i}$ (Stage~2 only)
      \STATE \quad Fit model on $\mathcal{D}_{\text{train}}$
      \STATE \quad Predict at patient $i$'s observation times
      \STATE \quad Record per-patient metrics
    \ENDFOR
    \STATE Aggregate metrics across $N$ folds
    \STATE Compute .632+ bootstrap CIs
  \end{algorithmic}
\end{algorithm}

All model fitting, feature transformation (PCA), and severity estimation
are performed within each fold. This ensures that LOPO-CV metrics reflect
true out-of-sample performance without information leakage.

For each held-out patient $i$, predictions are evaluated at the patient's
observed timepoints. For Stage~2, the held-out patient's severity is
re-estimated from their observations using the population model fitted on the
remaining patients (\cref{eq:severity_posterior}).

% ---------------------------------------------------------------------------
% 8.6 Volume Decoding and Uncertainty
% ---------------------------------------------------------------------------
\subsection{Volume Representation and Uncertainty}
\label{sec:growth:decoder}

\subsubsection{Volume Representation}

All models in Stages~1 and~2 operate on log-transformed whole-tumour volume:
\begin{equation}
  v = \log(V_{\text{WT}} + 1),
  \label{eq:log_vol}
\end{equation}
where $V_{\text{WT}}$ is the whole-tumour volume in mm$^3$. This is the
single clinically relevant growth endpoint for meningiomas. The
log-transformation stabilises variance and makes growth rates approximately
additive.

For Stage~3, the SDP semantic head decodes latent predictions to
log-volumes:
\begin{equation}
  \hat{v} = \pi_{\text{vol}}(\hat{\Vol})
    = \mathbf{W}_v \hat{\Vol} + \mathbf{b}_v,
  \quad \hat{v} \in \R^1,
  \label{eq:vol_decode}
\end{equation}
with de-normalisation $V_{\text{pred}} = \exp(\hat{v} \cdot \sigma_{\text{vol}}
+ \mu_{\text{vol}}) - 1$.

\subsubsection{Uncertainty Propagation}

For GP models that provide predictive variance $\sigma^2_{\hat{v}}$ in
log-volume space, the physical volume uncertainty is propagated through the
exponential via the delta method:
\begin{equation}
  \sigma_{V_{\text{pred}}} \approx V_{\text{pred}} \cdot \sigma_{\hat{v}}
    \cdot \sigma_{\text{vol}},
  \label{eq:delta_method}
\end{equation}
where the multiplicative structure reflects the log-normal nature of volume
uncertainty. For prediction intervals, we report the 95\% credible interval
from the GP posterior.

% ---------------------------------------------------------------------------
% 8.7 Failure Recovery
% ---------------------------------------------------------------------------
\subsection{Failure Recovery Protocol}
\label{sec:growth:failure}

If GP hyperparameter fitting diverges or produces degenerate results, the
following recovery steps are applied in order:
\begin{enumerate}[leftmargin=*]
  \item Add diagonal jitter to the kernel matrix:
    $\mathbf{K} \leftarrow \mathbf{K} + 10^{-6}\,\mathbf{I}$.
  \item Constrain length-scale bounds: $\ell \in [1, 120]$ months.
  \item Increase L-BFGS-B random restarts from 5 to 10.
  \item Fall back to a simpler kernel (Mat\'{e}rn-5/2 $\to$ SE).
  \item If all fail: report per-dimension diagnostics and halt.
\end{enumerate}

For Stage~2 severity estimation, if the joint optimisation fails to converge:
\begin{enumerate}[leftmargin=*]
  \item Reduce the growth function complexity (CMNN $\to$ sigmoid).
  \item Initialise severities from the empirical growth rate ranking.
  \item Fix population parameters and estimate severities only (profile
    likelihood).
\end{enumerate}
